\documentclass{article}
\usepackage[margin=1in]{geometry}
\usepackage{float}
\usepackage{cite}
\usepackage{amsmath,amssymb,amsfonts}
\usepackage{multirow}
\usepackage{graphicx}
\usepackage{booktabs}
\usepackage{textcomp}
\usepackage{hyperref}
\usepackage{xcolor}
\usepackage{array}
\usepackage{caption}

\def\BibTeX{{\rm B\kern-.05em{\sc i\kern-.025em b}\kern-.08em
    T\kern-.1667em\lower.7ex\hbox{E}\kern-.125emX}}

\begin{document}

\title{Wearable Accelerometry for Gait Assessment and Six-Minute Walk Distance Prediction in Pediatric-Onset Multiple Sclerosis: A Clinic and Home Comparison}

% \author{Arafat Rahman, Ariful Islam, Silvia Blemker, Rebecca Scharf, and Laura E. Barnes
% \thanks{Arafat Rahman, Ariful Islam, and Laura E. Barnes are with the
% Systems and Information Engineering Department, University of Virginia,
% Charlottesville, VA 22904 USA (e-mail: \{jgh6ds, xef4hb, lb3dp\}@virginia.edu).}
% \thanks{Silvia Blemker is with the Biomedical Engineering Department, University of Virginia.}
% \thanks{Rebecca Scharf is with the Department of Pediatrics, University of Virginia School of Medicine.}}

\maketitle

\begin{abstract}
Walking difficulty is common in pediatric-onset multiple sclerosis (POMS), but it is unclear how much clinically relevant gait information can be recovered from unsupervised home accelerometry. We compared clinic and home accelerometry as markers of walking impairment and as predictors of six-minute walk distance (6MWD) in 38 participants with POMS and 63 healthy controls. Participants completed a supervised six-minute walk test (6MWT) and seven days of hip-worn accelerometry, from which we derived 51 clinic gait and sway features and 151 home bout-level and whole-recording activity features. Compared with controls, participants with POMS showed lower cadence, lower walking intensity, and greater mediolateral sway normalized to walking effort. In the clinic, intensity-related features correlated most strongly with 6MWD. Across subjects, the strongest agreement between clinic and home measures occurred when home features were summarized at the 90th percentile of walking bouts, suggesting that the 6MWT reflects near-peak rather than habitual free-living performance. Continuous-wavelet spectral features and stride-timing variability were associated with fatigue scores.
For 6MWD prediction, Ridge regression achieved $R^2 = 0.81$ (MAE $= 31.2$ m) from clinic features and $R^2 = 0.45$ (MAE $= 56.0$ m) from a fully clinic-free home pipeline. Clinic prediction performed better in POMS than in healthy controls ($R^2 = 0.87$, MAE $= 29.0$ m vs $R^2 = 0.64$, MAE $= 32.5$ m), whereas home prediction was weaker in both groups ($R^2 = 0.38$ vs $0.21$). A compact 24-feature model retained much of the clinic signal ($R^2 = 0.69$), and the roughly 25 m MAE gap between clinic and home persisted across seven model families, indicating that the main limitation is the measurement setting rather than the regression algorithm. Home accelerometry therefore provides a useful but attenuated surrogate for clinic 6MWT performance in POMS.
\end{abstract}

\pagenumbering{arabic}

%% ============================================================
\section{Introduction}\label{sec:intro}
%% ============================================================

Multiple sclerosis (MS) is a chronic demyelinating disease of the central nervous system that affects about 2.8 million people worldwide~\cite{compston2008multiple, dobson2019multiple}. Between 50 and 80 percent of people with MS develop walking problems over time, which reduce daily function and quality of life~\cite{larocca2011impact}. These problems include slower walking speed, lower cadence, shorter steps, and poor balance~\cite{sosnoff2012quantifying, kalron2015association, hobart2003measuring}. Fatigue is the most commonly reported MS symptom~\cite{freal1984symptomatic, krupp1989fatigue} and makes walking worse over time~\cite{confavreux2000relapses}.

Pediatric-onset MS (POMS) is defined as symptom onset before age 18 and accounts for 3 to 10 percent of MS cases~\cite{chitnis2009demographics, banwell2007pediatric, yeh2009pediatric}. Compared to adults, POMS patients have higher relapse rates~\cite{gorman2009increased} and often reach disability milestones at a younger age~\cite{renoux2007natural, mckay2019long}. Even when overall disability is low, POMS participants are less active than their peers and show subtle walking deficits~\cite{brenton2022sixminute}. Despite these concerns, most MS gait studies are done in adult patients, and little is known about POMS gait when measured in a free-living setting.

The six-minute walk test (6MWT) is a standard measure of walking capacity in MS and in POMS~\cite{goldman2008evaluation, ats2002guidelines, brenton2022sixminute}. It gives a single number (six-minute walk distance, 6MWD) that summarizes cardiopulmonary, musculoskeletal, and neurological function. Clinic testing is accurate but gives only periodic snapshots of walking function and does not reflect day-to-day activity at home. Wearable accelerometers can fill this gap by continuously recording movement during daily life~\cite{motl2017wearable, alexander2021wearable, karas2023systematic}. Hip or waist-worn accelerometers can be used to measure cadence, harmonic ratio, stride regularity, and postural sway~\cite{sosnoff2012quantifying, igarashi2024accelerometer, moon2017monitoring, mancini2012isway, meyer2024assessing, meyer2023chest}. Activity-based measures from wearables have been linked to walking capacity in adult MS~\cite{motl2011accelerometry, carpinella2023automatic, silva2024comprehensive, fitzgerald2026activity, bourjeily2025accelerometry} and to disease progression in other neurological conditions~\cite{gupta2023athome}.

Wearable gait studies in POMS are rare. The only study that combined accelerometry with the 6MWT in POMS reported group-level activity differences but did not test predictive models or cross-setting comparisons~\cite{brenton2022sixminute}. Several open questions remain. First, how strongly do clinic-based and home-based features differ between POMS and healthy controls? Second, which features carry the most information about 6MWD? Third, do clinic features and home features measure the same underlying walking trait, or do they capture different things? Fourth, do features correlate with fatigue and disability scores, and in which setting? Fifth, can a home-based model approach the accuracy of a clinic-based model for predicting 6MWD?

We address these five questions in this paper. Specifically, we ask:

\begin{enumerate}
    \item \textbf{RQ1}: Which wearable features, measured in the clinic 6MWT and in free-living home monitoring, differ between POMS and healthy controls, and with what effect size?
    \item \textbf{RQ2}: Which clinic and home features are most correlated with measured 6MWD, and does the pattern hold within POMS alone?
    \item \textbf{RQ3}: How strongly do matched clinic and home features correlate across the same subject, and which summary statistic of the home distribution best recovers the clinic measurement?
    \item \textbf{RQ4}: Within POMS, how do clinic and home features correlate with EDSS disability, MFIS fatigue (total, physical, cognitive, psychosocial), and BDI depression?
    \item \textbf{RQ5}: Can a seven-day home accelerometry model predict 6MWD with accuracy that approaches a clinic-based 6MWT model, for the full cohort and within POMS alone?
\end{enumerate}

Our overarching question is whether seven-day home accelerometry can approach clinic 6MWT utility for describing walking capacity and clinical burden in POMS. The answers we report combine group comparisons, correlation analyses, and prediction models, giving a complete picture of what home monitoring can and cannot do in this population.

%% ============================================================
\section{Related Work}\label{sec:related}
%% ============================================================

\subsection{Wearable Gait Assessment in MS}

Many studies have used wearable sensors to measure gait in adult MS and related conditions. Motl et al.~\cite{motl2017wearable} and Alexander et al.~\cite{alexander2021wearable} reviewed wearable approaches that complement episodic clinic assessments of walking in MS. Sosnoff et al.~\cite{sosnoff2012quantifying} showed, using an instrumented walkway, that people with MS often walk more slowly and with greater gait variability than controls. In subacute stroke, Igarashi et al.~\cite{igarashi2024accelerometer} reported that trunk-acceleration metrics including normalized RMS measures could discriminate gait-independent from gait-dependent patients. In cerebellar ataxia, Phan et al.~\cite{phan2020random} used random forests on wearable gait features and showed strong agreement with expert-rated gait severity. Postural sway has been quantified from body-worn sensors in controlled balance tasks~\cite{mancini2012isway} and adapted for chest-worn assessment in MS~\cite{meyer2024assessing, meyer2023chest}.

\subsection{Walking Capacity and Home Monitoring}

The 6MWT is widely used in MS~\cite{goldman2008evaluation, ats2002guidelines}. Brenton et al.~\cite{brenton2022sixminute} validated the 6MWT in POMS and showed that POMS participants walk shorter distances and cannot accelerate in the last minute of the test. For free-living monitoring, Gashi et al.~\cite{gashi2024modeling} separated MS from controls with $F_1 = 0.82$ using wearable features. Bou Rjeily et al.~\cite{bourjeily2025accelerometry} showed that home accelerometry can detect MS disability status. Fitzgerald et al.~\cite{fitzgerald2026activity} showed that decreases in daytime activity track brain atrophy and disability progression in longitudinal follow-up. Carpinella et al.~\cite{carpinella2023automatic} estimated the two-minute walk distance from wearables in adult MS, reaching $r = 0.71$ in the clinic and $r = 0.58$ at home. Straczkiewicz et al.~\cite{straczkiewicz2023walking} describe an approach to walking-bout detection from wearable accelerometers. Gupta et al.~\cite{gupta2023athome} showed that home wearable features can track disease progression in amyotrophic lateral sclerosis with higher sensitivity than clinical scales.

\subsection{Gap in the Literature}

Two gaps motivate this work. First, most wearable MS studies are based on adult cohorts and rarely include pediatric-onset patients. POMS participants have lower disability but subtle gait changes, and the features that work for adult MS may not transfer directly. Second, few studies compare clinic and home features within the same subjects in a systematic way. Without that comparison, it is not clear whether a home-based assessment can substitute for a supervised clinic test. Our paper fills both gaps by using a pediatric-only cohort and by comparing clinic and home features at the subject level.

%% ============================================================
\section{Methods}\label{sec:methods}
%% ============================================================

\begin{figure}[!ht]
    \centering
    \includegraphics[width=0.99\linewidth]{figures/fig_overview_diagram.png}
    \caption{Study overview. Each of 101 participants completed a supervised six-minute walk test (6MWT) in the clinic and wore a hip-worn accelerometer for seven days at home. From the clinic recording we computed 51 gait features. From the home recording we detected walking bouts and computed 151 per-bout gait and activity features. The same subjects contribute to both the clinic and home analyses, so we can compare features, associations, and predictions directly.}
    \label{fig:overview}
\end{figure}

\subsection{Participants}\label{sec:participants}

Participants with POMS were recruited from the University of Virginia (UVA) Pediatric MS Clinic between 2016 and 2024. All recruited participants met the 2010 revised McDonald criteria for relapsing-remitting MS with onset before age 18~\cite{polman2011diagnostic}. Age-matched healthy controls were recruited from the local community and UVA general pediatric clinics. The final cohort was 101 participants (38 POMS, 63 healthy controls) after excluding two POMS subjects for data-quality reasons. Table~\ref{tab:demographics} gives demographic and clinical characteristics. The POMS group was slightly younger ($16.5 \pm 2.6$ vs.\ $18.0 \pm 3.1$ years, $p = 0.007$), had higher BMI ($27.0 \pm 6.1$ vs.\ $22.9 \pm 3.9$, $p < 0.001$), and walked a shorter 6MWD ($576.1 \pm 99.1$ vs.\ $674.7 \pm 81.4$ m, $p < 0.001$). POMS participants reported higher physical fatigue (MFIS Physical $8.9 \pm 7.5$ vs.\ $4.7 \pm 4.9$, $p = 0.009$) and higher depressive symptoms (BDI $8.0 \pm 8.4$ vs.\ $4.2 \pm 5.8$, $p = 0.004$). Median EDSS among POMS participants was $1.8 \pm 0.8$ and mean disease duration was $2.5 \pm 2.3$ years. Additional cohort details are reported by Brenton et al.~\cite{brenton2022sixminute}. The study was approved by the UVA Institutional Review Board. Written informed consent was obtained from all adult participants and written assent from minors with parental consent.

\begin{table}[!ht]
\centering
\caption{Demographic and clinical characteristics of the study cohort. Values are mean $\pm$ SD, with the range [min, max] shown in brackets for 6MWD. $p$-values from Mann-Whitney $U$ tests (continuous variables) or chi-squared test (sex). EDSS and disease duration are reported for POMS only. *$p<0.05$, **$p<0.01$, ***$p<0.001$.}
\label{tab:demographics}
\footnotesize
\begin{tabular}{@{}lcccc@{}}
\toprule
Variable & All ($n=101$) & POMS ($n=38$) & Healthy ($n=63$) & $p$-value \\
\midrule
Age (years) & $17.5 \pm 3.0$ & $16.5 \pm 2.6$ & $18.0 \pm 3.1$ & 0.007** \\
Sex (M/F) & 33/68 & 11/27 & 22/41 & 0.688 \\
Height (cm) & $168.6 \pm 10.2$ & $167.6 \pm 10.9$ & $169.2 \pm 9.8$ & 0.422 \\
Weight (kg) & $70.1 \pm 19.3$ & $76.7 \pm 23.8$ & $66.1 \pm 14.8$ & 0.027* \\
BMI (kg/m$^2$) & $24.4 \pm 5.2$ & $27.0 \pm 6.1$ & $22.9 \pm 3.9$ & $<$0.001*** \\
6MWD (m) & $637.6 \pm 100.5$ [339--1040] & $576.1 \pm 99.1$ [339--759] & $674.7 \pm 81.4$ [530--1040] & $<$0.001*** \\
\addlinespace
EDSS & --- & $1.8 \pm 0.8$ & --- & --- \\
Disease Duration (yrs) & --- & $2.5 \pm 2.3$ & --- & --- \\
\addlinespace
MFIS Total & $15.0 \pm 13.9$ & $19.9 \pm 16.6$ & $12.0 \pm 11.1$ & 0.027* \\
MFIS Physical & $6.3 \pm 6.3$ & $8.9 \pm 7.5$ & $4.7 \pm 4.9$ & 0.009** \\
MFIS Cognitive & $7.6 \pm 7.5$ & $9.8 \pm 8.9$ & $6.3 \pm 6.2$ & 0.141 \\
MFIS Psychosocial & $1.0 \pm 1.2$ & $1.3 \pm 1.4$ & $0.9 \pm 1.1$ & 0.202 \\
BDI & $5.6 \pm 7.1$ & $8.0 \pm 8.4$ & $4.2 \pm 5.8$ & 0.004** \\
\bottomrule
\end{tabular}
\end{table}

\subsection{Data Collection}\label{sec:collection}

\subsubsection{Clinic 6MWT}
Each participant walked for six minutes in a standardized hospital corridor while wearing an ActiGraph GT3X hip-mounted accelerometer at 30 Hz. Total distance was recorded in feet and converted to meters ($\times 0.3048$). On the same visit, the EDSS~\cite{kurtzke1983rating}, MFIS~\cite{krupp1989fatigue}, and BDI~\cite{beck1961inventory} were administered.

\subsubsection{Home Monitoring}
After the clinic visit, participants wore a hip-mounted ActiGraph accelerometer for seven days during all waking activities, except during water-based activities. The device sampled continuously at 30 Hz. Daily wear time, estimated from minute-level vector-magnitude variability using a standard accelerometer wear/nonwear detection rule, was comparable between POMS and healthy controls (median 59 vs.\ 53 hours; Mann-Whitney $U$, $p = 0.19$; Supplementary Figure~\ref{fig:supp_wear_median_bout}), indicating that between-cohort compliance was not a confounding factor for subsequent home analyses.

\subsection{Data Preprocessing}\label{sec:preprocessing}

\subsubsection{Clinic Data}
Raw tri-axial acceleration from the 6MWT was processed in four steps: (1) gravity was estimated with a 0.25 Hz fourth-order Butterworth lowpass filter and subtracted; (2) the vertical axis was aligned with the estimated gravity vector using Rodrigues rotation; (3) principal component analysis on the horizontal plane resolved the yaw, giving anteroposterior (AP), mediolateral (ML), and vertical (VT) axes; (4) a bandpass filter at 0.25 to 2.5 Hz isolated the gait band.

\subsubsection{Home Data}
Home data were processed per walking bout. Walking bouts were detected with a three-stage pipeline: per-second Euclidean Norm Minus One (ENMO) thresholding at $\geq 0.015\,g$ with a minimum 10 s duration, harmonic ratio verification in 10 s sliding windows, and merging of adjacent bouts separated by gaps of at most 5 s. The 10 s minimum-duration floor was chosen empirically and was confirmed to be the best operating point for 6MWD prediction in a sensitivity analysis (Supplementary Figure~\ref{fig:supp_bout_threshold}): raising the threshold to 30--240\,s reduced the number of retained bouts and degraded the home Bout+Act-Top20+Demo MAE from 56.0\,m to 62--66\,m. All detected bouts were kept without further quality filtering. Each bout was then processed the same way as the clinic data: gravity removal, axis alignment, PCA-based yaw, and bandpass filtering.

\subsection{Feature Extraction}\label{sec:features}

Table~\ref{tab:features} lists all features by category and setting. The clinic side has 51 features in three families. The home side has 151 features that summarize per-bout gait and whole-recording activity.

\subsubsection{Clinic Gait Features (11)}
These features capture the temporal structure and intensity of walking:
\begin{itemize}
    \item \textbf{cadence\_hz}: walking cadence in Hz, estimated from the peak of the vertical power spectral density in the 0.5 to 3.5 Hz gait band. Reduced cadence is a well-established sign of gait impairment in MS~\cite{sosnoff2012quantifying}.
    \item \textbf{step\_time\_cv\_pct}: coefficient of variation of step times in percent, computed from peak-to-peak intervals in the vertical bandpass signal. Higher step-time variability is associated with greater disability and less regular gait in MS~\cite{kalron2015association}.
    \item \textbf{acf\_step\_regularity}: autocorrelation of the vertical bandpass signal at the step lag, used as a step-regularity index; values close to 1 indicate more repeatable stepping.
    \item \textbf{hr\_ap, hr\_vt}: harmonic ratios in the anteroposterior and vertical directions, computed as the ratio of even to odd harmonics of the gait frequency. Higher values reflect smoother and more symmetric gait~\cite{igarashi2024accelerometer}.
    \item \textbf{ml\_rms\_g}: root mean square of mediolateral dynamic acceleration in gravitational units. ML acceleration reflects the magnitude of side-to-side trunk oscillation and has been shown to discriminate gait independence in neurological populations~\cite{igarashi2024accelerometer}.
    \item \textbf{ml\_spectral\_entropy}: normalized spectral entropy of the ML power spectrum in the cadence-relative frequency band. Higher entropy reflects a less concentrated spectral profile.
    \item \textbf{jerk\_mean\_abs\_gps}: mean absolute jerk of the 3D velocity magnitude in $g$/s. Jerk quantifies the smoothness of movement and is used in postural-control assessment~\cite{mancini2012isway}.
    \item \textbf{enmo\_mean\_g}: mean Euclidean Norm Minus One (ENMO), computed as $\max(\sqrt{a_x^2 + a_y^2 + a_z^2} - 1\,g,\; 0)$. ENMO is a standard measure of movement intensity that removes the gravitational component~\cite{vanhees2013separating}.
    \item \textbf{cadence\_slope\_per\_min}: linear slope of cadence over the six one-minute segments of the 6MWT, summarizing within-test change in cadence that may reflect deceleration over the walk~\cite{qureshi2016relationship}.
    \item \textbf{vt\_rms\_g}: root mean square of vertical dynamic acceleration in $g$. VT RMS captures the vertical oscillation amplitude during the gait cycle and has been used as a gait-quality marker in free-living monitoring~\cite{storm2020free, igarashi2024accelerometer}.
\end{itemize}

\subsubsection{Clinic CWT Features (28)}
We applied a continuous wavelet transform (CWT)~\cite{torrence1998practical} with a Morlet wavelet to the vector magnitude acceleration. The CWT gives a time-frequency representation that exposes how spectral content of walking evolves over the test. For each of the six one-minute segments we computed ten base features:
\begin{itemize}
    \item \textbf{mean\_energy}: average wavelet power across all time and frequency bins.
    \item \textbf{high\_freq\_energy}: mean wavelet power above 3.5 Hz, capturing high-frequency gait dynamics such as foot-strike impacts.
    \item \textbf{dominant\_freq}: frequency with the highest time-averaged wavelet power.
    \item \textbf{estimated\_cadence}: cadence (steps/min) estimated from the CWT gait band (0.5 to 3.5 Hz).
    \item \textbf{max\_power\_freq}: frequency of the global maximum in the time-frequency plane.
    \item \textbf{freq\_variability}: standard deviation of the dominant frequency across one-second windows; higher cadence-related frequency variability may indicate less stable gait rhythm.
    \item \textbf{freq\_cv}: coefficient of variation of the dominant frequency.
    \item \textbf{wavelet\_entropy}: Shannon entropy of the normalized wavelet power distribution, measuring the uniformity of power across frequencies.
    \item \textbf{fundamental\_freq}: fundamental gait frequency estimated from the FFT of the segment.
    \item \textbf{harmonic\_ratio}: ratio of even to odd FFT harmonic amplitudes, reflecting gait symmetry and smoothness in the frequency domain.
\end{itemize}
The ten base features were aggregated across the six segments as mean and standard deviation (20 features), plus a linear slope and slope $R$ value for four fatigue-sensitive descriptors (mean\_energy, high\_freq\_energy, freq\_variability, wavelet\_entropy), giving 28 features total. Slope features capture gradual changes over the walk that are thought to reflect motor fatigue~\cite{qureshi2016relationship}. The CWT feature set intentionally provides segment-level, wavelet-based estimators of cadence (fundamental\_freq, dominant\_freq, estimated\_cadence) that complement the Welch-PSD cadence computed over the full walk in the Gait feature set; the two estimators are correlated ($r \approx 0.72$) but not identical, reflecting their different time-frequency resolutions, and Ridge regression with feature standardization handles the residual correlation without overfitting.

\subsubsection{Clinic WalkSway Features (12)}
Raw sway metrics such as ML range, path length, and jerk scale with walking speed: faster walkers generate larger accelerations in all directions, which inflates sway measures regardless of actual postural instability~\cite{mancini2012isway, meyer2023chest}. To remove this confound, each sway metric was divided by ENMO (or by ENMO$^2$ where area scales quadratically). The 12 features are:
\begin{itemize}
    \item \textbf{ml\_range\_norm}: peak-to-peak ML range divided by ENMO, representing lateral excursion per unit of walking effort.
    \item \textbf{ml\_path\_length\_norm}: cumulative absolute ML displacement divided by ENMO.
    \item \textbf{ml\_jerk\_rms\_norm}: RMS of ML jerk (second derivative) divided by ENMO, measuring smoothness of lateral trunk control.
    \item \textbf{ap\_rms\_norm, ap\_range\_norm}: AP RMS and peak-to-peak range divided by ENMO.
    \item \textbf{sway\_ellipse\_norm}: area of the 95\% confidence ellipse in the ML-AP plane divided by ENMO$^2$, a standard measure of overall sway volume in postural-control studies~\cite{mancini2012isway, meyer2023chest}.
    \item \textbf{ml\_velocity\_rms\_norm}: RMS of ML velocity (first derivative) divided by ENMO.
    \item \textbf{stride\_ml\_cv}: coefficient of variation of stride-to-stride ML peak amplitudes. Stride-to-stride variability is a sensitive marker of gait instability in neurological populations~\cite{kalron2015association, phan2020random}.
    \item \textbf{ml\_ap\_ratio}: ML RMS divided by AP RMS, capturing the directional dominance of sway.
    \item \textbf{hr\_ml}: harmonic ratio in the ML direction, quantifying left-right gait symmetry.
    \item \textbf{ml\_over\_enmo}: ML RMS divided by ENMO, a simple normalized lateral-sway index.
    \item \textbf{ml\_over\_vt}: ML RMS divided by VT RMS, capturing lateral sway relative to vertical gait dynamics.
\end{itemize}

\subsubsection{Home Features (151)}
For each detected walking bout we computed 20 per-bout gait features (cadence, stride timing, step regularity, AP/ML/VT RMS, jerk, harmonic ratios, ENMO). Because subjects have variable bout counts, we aggregated per-bout features using six summary statistics per feature: median, interquartile range (IQR), 10th percentile, 90th percentile, maximum, and coefficient of variation (CV), giving 122 per-bout features in total. We also computed 29 whole-recording activity features spanning ENMO distribution moments and percentiles, proportions of sedentary, light, moderate, and vigorous activity, fragmentation indices (active-to-sedentary transition probability), and early/mid/late-day ENMO trends. Activity-intensity thresholds use the ENMO metric introduced by van Hees et al.~\cite{vanhees2013separating}. The home features give the free-living counterpart of the supervised clinic 6MWT at the subject level.

\subsubsection{Demographics (3 to 4)}
Demographic features were cohort (POMS vs.\ healthy), age, sex, and either height (clinic models) or BMI (home models).

\begin{table}[!ht]
\centering
\caption{Feature categories and counts. Clinic features come from the supervised 6MWT recording. Home features come from seven-day free-living accelerometry with automatic walking-bout detection.}
\label{tab:features}
\small
\begin{tabular}{@{}p{1.2cm}p{2.4cm}cp{8.5cm}@{}}
\toprule
Setting & Category & Count & Description \\
\midrule
Clinic & Gait & 11 & Cadence, step timing, autocorrelation, harmonic ratios, RMS, entropy, jerk, ENMO, cadence slope \\
\addlinespace
Clinic & CWT & 28 & 10 base wavelet descriptors $\times$ 3 aggregations (mean, SD, and a linear slope for four fatigue-sensitive features) \\
\addlinespace
Clinic & WalkSway & 12 & ENMO-normalized ML and AP sway metrics, sway ellipse, stride variability, harmonic ratios \\
\addlinespace
Home & Bout & 122 & 20 per-bout roots $\times$ 6 aggregations (median, IQR, p10, p90, max, CV) + 2 meta (total walk seconds, mean bout duration) \\
\addlinespace
Home & Activity & 29 & Whole-recording ENMO moments and percentiles, activity-intensity proportions, fragmentation, diurnal trends \\
\bottomrule
\end{tabular}
\end{table}

\subsection{Statistical Analysis}\label{sec:stats}

Group differences between POMS and healthy controls were tested with Mann-Whitney $U$ tests. Effect sizes were reported as Cohen's $d$ (positive $d$ means POMS is higher than healthy). $p$-values were adjusted for multiple testing across features within each setting using the Benjamini-Hochberg procedure~\cite{benjamini1995controlling}. Feature-to-outcome associations (6MWD, clinical scores) were quantified with Spearman rank correlations, which capture monotonic non-linear relationships.

\subsection{Prediction Models}\label{sec:models}

The main prediction model was Ridge regression~\cite{hoerl1970ridge}. Ridge was chosen because it handles correlated features well and is stable at small sample sizes, which matches our cohort of 101. The regularization parameter was $\alpha = 5$ for the clinic setting and $\alpha = 20$ for the home setting, both selected by prior experimentation. All models were evaluated using leave-one-subject-out (LOO) cross-validation. Features were standardized (zero mean, unit variance) within each training fold to prevent information leakage. For the clinic model, the full 55-feature set (Gait + CWT + WalkSway + Demo) was used without feature selection. For the home model, Spearman-based feature selection was run inside each LOO fold: the 20 per-bout features most correlated with 6MWD on the training set were kept, and the four demographic features were appended. We also compared Ridge against six alternative models (Lasso, ElasticNet, $k$-nearest neighbors, support vector regression, random forest, and XGBoost) under the same LOO protocol.

Model performance was summarized by $R^2$, mean absolute error (MAE) in meters, and Pearson $r$ between predicted and actual 6MWD. Ninety-five percent confidence intervals were obtained by bootstrapping the subject-level predictions with 2{,}000 resamples.

%% ============================================================
\section{Results}\label{sec:results}
%% ============================================================

The Results are organized to match the five research questions.

\subsection{RQ1: Features That Differ Between POMS and Healthy Controls}\label{sec:rq1}

Table~\ref{tab:mshealthy} lists the features that differed between groups at $p_{\mathrm{BH}} < 0.05$ in each setting, ranked by effect size. We kept features whose direction is consistent with gait literature (higher value means more pathology or less vigor). In the clinic, POMS participants showed lower spectral tempo (CWT fundamental frequency $d = -0.92$, estimated cadence $d = -0.82$, dominant frequency $d = -0.77$, cadence $d = -0.73$), lower walking intensity (ENMO $d = -0.83$, VT RMS $d = -0.78$), and lower absolute jerk ($d = -0.70$). Only one normalized sway feature reached significance with the expected direction: ML over ENMO was higher in POMS ($d = 0.66$, $p_{\mathrm{BH}} = 0.017$), meaning that for the same level of walking effort, POMS participants showed more lateral sway.

In the home setting, ten features were significantly different. Per-bout jerk ($d = -1.12$), AP RMS ($d = -1.02$), bout-duration coefficient of variation ($d = -0.98$), longest bout duration ($d = -0.77$), and 10th-percentile ENMO ($d = -0.75$) were all lower in POMS, indicating shorter, less vigorous, and less varied walking. Cross-bout AP RMS variability was higher in POMS ($d = 0.74$), reflecting more inconsistent gait from bout to bout. Activity-level features showed smaller effects: percentage of time in vigorous activity ($d = -0.33$) and 95th-percentile ENMO ($d = -0.27$) were lower in POMS. The pattern is consistent across settings: POMS participants walked less intensely and with more variability.

Figure~\ref{fig:bout_overview} shows the raw home walking-bout distribution. The violin panel shows that healthy controls have a wider spread of longest bout durations (with more subjects reaching multi-minute bouts), while POMS participants cluster at shorter durations. The scatter panel shows a positive relationship between number of bouts and longest bout duration in both groups. The survival panel confirms that the two distributions start together at short durations but separate past 60 s, where healthy controls keep more bouts of longer duration. Per-subject bout-duration box plots (Supplementary Figure~\ref{fig:supp_boxplots}) show the within-subject spread of bout durations and confirm that the between-cohort separation is visible at the individual level, not only in pooled statistics.

\begin{table}[!ht]
\centering
\caption{Features that differ significantly between POMS and healthy controls in each setting. POMS and Healthy columns show mean~$\pm$~SD in the native feature units. Rows are ordered by $|d|$ within each setting. Negative $d$ means POMS is lower than healthy. Significance levels use BH-adjusted $p$-values: *$p<0.05$, **$p<0.01$, ***$p<0.001$.}
\label{tab:mshealthy}
\footnotesize
\begin{tabular}{@{}llccrr@{}}
\toprule
Setting & Feature & POMS ($n=38$) & Healthy ($n=63$) & $d$ & $p_{\mathrm{BH}}$ \\
\midrule
Clinic (CWT)      & cwt\_fundamental\_freq\_mean  & $1.88 \pm 0.49$      & $2.19 \pm 0.20$      & $-0.92$ & 0.0007*** \\
Clinic (Gait)     & enmo\_mean\_g                 & $0.227 \pm 0.084$    & $0.304 \pm 0.098$    & $-0.83$ & 0.0009*** \\
Clinic (CWT)      & cwt\_estimated\_cadence\_mean & $138.2 \pm 36.8$     & $159.5 \pm 15.9$     & $-0.82$ & 0.0034** \\
Clinic (Gait)     & vt\_rms\_g                    & $0.454 \pm 0.147$    & $0.569 \pm 0.146$    & $-0.78$ & 0.0019** \\
Clinic (CWT)      & cwt\_dominant\_freq\_mean     & $2.30 \pm 0.61$      & $2.73 \pm 0.51$      & $-0.77$ & 0.0028** \\
Clinic (Gait)     & cadence\_hz                   & $2.04 \pm 0.19$      & $2.18 \pm 0.19$      & $-0.73$ & 0.0025** \\
Clinic (Gait)     & jerk\_mean\_abs\_gps          & $7.68 \pm 2.99$      & $9.90 \pm 3.23$      & $-0.70$ & 0.0039** \\
Clinic (WalkSway) & ml\_over\_enmo                & $1.16 \pm 0.34$      & $0.99 \pm 0.20$      & $+0.66$ & 0.0170* \\
\addlinespace
Home (Bout)       & g\_jerk\_mean\_med            & $1.38 \pm 0.73$      & $2.38 \pm 0.97$      & $-1.12$ & 0.0001*** \\
Home (Bout)       & g\_ap\_rms\_med               & $0.098 \pm 0.050$    & $0.157 \pm 0.061$    & $-1.02$ & 0.0001*** \\
Home (Bout)       & g\_duration\_sec\_cv          & $1.34 \pm 0.44$      & $2.06 \pm 0.86$      & $-0.98$ & 0.0001*** \\
Home (Bout)       & g\_duration\_sec\_max         & $606.5 \pm 415.1$    & $1163.2 \pm 853.0$   & $-0.77$ & 0.0007*** \\
Home (Bout)       & g\_enmo\_mean\_p10            & $0.0249 \pm 0.0035$  & $0.0284 \pm 0.0052$  & $-0.75$ & 0.0022** \\
Home (Bout)       & g\_ap\_rms\_cv                & $0.850 \pm 0.276$    & $0.682 \pm 0.188$    & $+0.74$ & 0.0009*** \\
Home (Bout)       & g\_acf\_step\_reg\_max        & $0.910 \pm 0.070$    & $0.936 \pm 0.037$    & $-0.49$ & 0.0102* \\
Home (Act)        & act\_pct\_vigorous            & $0.070 \pm 0.048$    & $0.085 \pm 0.042$    & $-0.33$ & 0.0395* \\
Home (Act)        & act\_enmo\_p95                & $0.123 \pm 0.064$    & $0.137 \pm 0.044$    & $-0.27$ & 0.0280* \\
\bottomrule
\end{tabular}
\end{table}

\begin{figure}[!ht]
    \centering
    \includegraphics[width=0.99\linewidth]{figures/bout_distribution_overview.png}
    \caption{Home walking-bout characteristics by cohort ($n=101$). Left: violin plot of the longest sustained bout per subject. Middle: violin plot of the per-subject 90th percentile of bout intensity (mean ENMO, $g$). Right: pooled bout-duration survival curve with 95\% bootstrap CI. POMS subjects cluster at shorter bout durations and have fewer long sustained bouts.}
    \label{fig:bout_overview}
\end{figure}

\subsection{RQ2: Features Correlated with 6MWD}\label{sec:rq2}

Figure~\ref{fig:corr_6mwd} shows the Spearman correlation between each top-ranked feature and 6MWD in all 101 subjects, POMS only, and healthy only. In the clinic panel, gait intensity and tempo features had the strongest correlations: ENMO ($\rho = 0.80$), VT RMS ($\rho = 0.76$), mean absolute jerk ($\rho = 0.73$), CWT fundamental frequency ($\rho = 0.68$), and cadence ($\rho = 0.68$). All are positive, meaning higher vigor and higher tempo are associated with longer 6MWD. ML over ENMO is the only pathology ratio in the top 10 and correlates negatively ($\rho = -0.42$), meaning more relative lateral sway is associated with shorter 6MWD. The pattern holds within the POMS subgroup, with most correlations in the range $\rho = 0.58$ to $0.79$.

In the home panel, correlations with 6MWD are smaller in magnitude but directionally consistent. Longer maximum bout duration, higher percentage of vigorous activity, more varied bout durations, and higher peak activity intensity all correlate positively with 6MWD ($\rho = 0.34$ to $0.37$). Stride-to-bout AP RMS variability correlates negatively ($\rho = -0.35$). Free-living correlations are smaller because walking bouts at home are shorter and more heterogeneous than the supervised 6MWT, but the direction of each correlation matches the clinic setting.

\begin{figure}[!ht]
    \centering
    \includegraphics[width=0.99\linewidth]{figures/heatmap_feature_6mwd_corr_top10.png}
    \caption{Spearman correlations between accelerometry features and 6MWD in the full cohort (All), the POMS subgroup, and the healthy subgroup. Clinic features (left) and home features (right) are shown side by side. Gait intensity and tempo features dominate the clinic side. Home correlations are smaller but directionally consistent. Significance stars use raw $p$-values: *$p<0.05$, **$p<0.01$, ***$p<0.001$.}
    \label{fig:corr_6mwd}
\end{figure}

\subsection{RQ3: Correlation Between Clinic and Home Features}\label{sec:rq3}

If the supervised 6MWT and free-living walking measure the same underlying trait, we expect clinic features and home features to correlate across subjects. To test this, we paired each of the top five clinic Gait features with its home counterpart (the same variable computed in free-living walking bouts) and computed Spearman correlations for each of six home aggregation statistics (median, IQR, 10th and 90th percentile, maximum, and CV). We also computed a $5 \times 5$ cross-correlation panel to see whether features spill into each other.

Figure~\ref{fig:clinic_home_corr} shows the result. In the full cohort, clinic ENMO is most correlated with the 90th-percentile home ENMO ($\rho = 0.33$), clinic jerk with the 90th-percentile home jerk ($\rho = 0.31$), and so on. The 90th-percentile column is consistently the densest cluster of significant positive correlations. Clinic and home tempo (cadence) features agree at both the median ($\rho = 0.30$) and the 90th percentile ($\rho = 0.30$) of home cadence, consistent with cadence being driven by habit rather than by effort.

This pattern supports a simple interpretation: the six-minute walk test pushes each subject toward the upper end of their own free-living effort distribution. The clinic measurement lines up most cleanly with the subject's 90th-percentile home value, not with their median or maximum. In the POMS subgroup the correlations are similar in direction but smaller in magnitude ($\rho = 0.24$ to $0.37$), reflecting the smaller sample size. In the healthy subgroup the correlations are weaker still because healthy subjects have less between-subject variation in both settings.

\begin{figure}[!ht]
    \centering
    \includegraphics[width=0.99\linewidth]{figures/heatmap_clinic_home_feature_corr.png}
    \caption{Clinic-to-home correlation for five matched Gait features. Rows are the clinic feature. Columns are the 90th-percentile home counterpart. Each panel shows one cohort group (All, POMS, Healthy). Diagonal cells are the direct matches; off-diagonal cells show cross-feature spillover. The densest cluster of significant cells lies at the 90th-percentile column, meaning the supervised 6MWT tracks the upper decile of free-living walking rather than the median.}
    \label{fig:clinic_home_corr}
\end{figure}

\subsection{RQ4: Correlations with Clinical Scores}\label{sec:rq4}

Within the POMS subgroup ($n = 38$) we tested how each feature correlates with EDSS, the four MFIS subscales (Total, Physical, Cognitive, Psychosocial), and BDI. Figures~\ref{fig:clin_corr_clinic} and~\ref{fig:clin_corr_home} show the top 10 features per setting, restricted to features whose sign matches the expected direction from the gait literature.

In the clinic setting, the strongest associations come from CWT spectral features. Mean CWT high-frequency energy correlates positively with MFIS Psychosocial ($\rho = 0.44$), MFIS Physical ($\rho = 0.39$), and MFIS Total ($\rho = 0.33$), meaning higher high-frequency content is associated with more fatigue. The standard deviation of the harmonic ratio correlates with MFIS Physical ($\rho = 0.40$) and MFIS Cognitive ($\rho = 0.32$), meaning larger session-to-session changes in gait symmetry go with more fatigue. Other significant rows include standard deviation of CWT high-frequency energy, wavelet entropy mean, slope of high-frequency energy, standard deviation of frequency variability, AP harmonic ratio, standard deviation of fundamental frequency, and the normalized ML path length. Most associations are modest ($\rho$ in the 0.2 to 0.4 range), consistent with the small POMS sample.

In the home setting, correlations are stronger overall and cluster around two axes. Shorter and less varied bouts go with worse scores: maximum bout duration correlates negatively with MFIS Cognitive ($\rho = -0.45$) and MFIS Total ($\rho = -0.39$). Bout-duration CV and activity-to-sedentary transition probability show the same pattern. Slower and more variable stride timing go with worse scores: mean stride-time at the 90th percentile correlates with MFIS Total ($\rho = 0.45$), MFIS Cognitive ($\rho = 0.44$), and MFIS Physical ($\rho = 0.39$). Several stride-timing variability features (90th-percentile CV, median CV, 90th-percentile SD, IQR of SD) all correlate positively with MFIS Total at $\rho = 0.40$ to $0.42$. BDI tracks the same features. EDSS has weaker associations in this sample, possibly because EDSS values cluster at the lower end of the scale (median $1.8 \pm 0.8$) and give little between-subject variation.

Taken together, RQ4 shows that home features give a more complete picture of fatigue and depression burden than clinic features alone. Two interpretable axes explain most of the signal: bout length and stride-timing stability. Both are measurable at home with a single hip-worn accelerometer.

\begin{figure}[!ht]
    \centering
    \includegraphics[width=0.99\linewidth]{figures/heatmap_clinical_corr_clinic.png}
    \caption{Spearman correlations between clinic accelerometry features and clinical scores in the POMS subgroup ($n=38$). Rows are the top 10 features whose sign matches gait-literature expectations. CWT irregularity and high-frequency energy rise with fatigue (MFIS) and depression (BDI) scores.}
    \label{fig:clin_corr_clinic}
\end{figure}

\begin{figure}[!ht]
    \centering
    \includegraphics[width=0.99\linewidth]{figures/heatmap_clinical_corr_home.png}
    \caption{Spearman correlations between home accelerometry features and clinical scores in the POMS subgroup ($n=38$). Two axes explain most of the signal: shorter and less varied walking bouts, and slower and more variable stride timing.}
    \label{fig:clin_corr_home}
\end{figure}

\subsection{RQ5: 6MWD Prediction Performance}\label{sec:rq5}
Figure~\ref{fig:model_bar} compares seven models across both settings. Linear models (Ridge, ElasticNet, Lasso) give the lowest clinic MAE (30 to 33 m), and the non-linear models (XGBoost, RF, KNN) are all worse (43 to 56 m). In the home setting the spread is smaller: SVR and Lasso are lowest (55 m) and KNN is highest (66 m), with Ridge in the middle at 56 m. The clinic-to-home gap is about 25 m of MAE across all models, which indicates that the gap is driven by the setting, not by the model choice.

\begin{figure}[!ht]
    \centering
    \includegraphics[width=0.99\linewidth]{figures/fig_model_comparison_mae.png}
    \caption{6MWD prediction MAE by model for clinic and home settings. Error bars are the bootstrap-derived standard-deviation equivalents of the 95\% confidence intervals. Linear models dominate in the clinic; the home setting is flatter across models, with a roughly 25 m MAE gap from clinic to home that is consistent across all seven models.}
    \label{fig:model_bar}
\end{figure}


Table~\ref{tab:results} reports Ridge regression prediction performance for each feature set across both settings. In the clinic, the full Gait + CWT + WalkSway + Demo feature set (55 features) gave the best result: $R^2 = 0.81$ (95\% CI $[0.71, 0.88]$), MAE $= 31.2$ m ($[25.5, 37.1]$), and $r = 0.90$ ($[0.85, 0.94]$). The gait-only model reached $R^2 = 0.68$ and the combined model gained an extra 13 points by adding CWT, WalkSway, and demographics. CWT and WalkSway alone reached only $R^2 = 0.36$ and $0.40$, confirming that gait intensity and tempo carry most of the clinic signal while CWT and WalkSway add complementary information.

In the home setting, the best clinic-free model was Bout + Act-Top20 + Demo (24 features), with $R^2 = 0.45$ ($[0.30, 0.55]$), MAE $= 56.0$ m ($[46.8, 65.5]$), and $r = 0.67$ ($[0.57, 0.76]$). Demographics alone reached $R^2 = 0.36$, and the home feature model without demographics reached $R^2 = 0.18$. The combination of per-bout home features and self-reported demographics is therefore essential. Applying the 51 clinic-style features directly to home data gave only $R^2 = 0.28$, suggesting that clinic feature definitions do not transfer cleanly to free-living walking. Running the Bout+Act-Top20+Demo pipeline on the clinic 6MWT recording also gave a competitive $R^2 = 0.69$ (MAE $= 39.6$ m), showing that the per-bout aggregation strategy works in both settings and is close to the full 55-feature clinic model with less than half the features.

\begin{table}[!ht]
\centering
\caption{6MWD prediction performance by feature set and setting. Ridge regression with leave-one-out cross-validation. Values are mean [95\% bootstrap CI] from 2{,}000 resamples. Best per setting is in \textbf{bold}.}
\label{tab:results}
\small
\begin{tabular}{@{}lcccccc@{}}
\toprule
Feature Set & \#f & Clinic $R^2$ & Clinic MAE (m) & Clinic $r$ & Home $R^2$ & Home MAE (m) \\
\midrule
Gait                 & 11 & 0.68 & 42.7 & 0.83 & 0.15 & 70.1 \\
CWT                  & 28 & 0.36 & 60.2 & 0.60 & 0.15 & 67.9 \\
WalkSway             & 12 & 0.40 & 54.2 & 0.63 & 0.06 & 73.3 \\
Demo                 & 4  & 0.36 & 60.8 & 0.60 & 0.36 & 60.8 \\
Bout+Act-Top20       & 20 & 0.63 & 44.3 & 0.79 & 0.18 & 67.2 \\
Bout+Act-Top20+Demo  & 24 & 0.69 & 39.6 & 0.83 & \textbf{0.45} & \textbf{56.0} \\
Gait+CWT+WS+Demo     & 55 & \textbf{0.81} & \textbf{31.2} & \textbf{0.90} & 0.28 & 63.8 \\
\bottomrule
\end{tabular}
\end{table}

Table~\ref{tab:error} reports per-cohort error for the two best models. The clinic model performs better in POMS ($R^2 = 0.87$, MAE $= 29.0$ m, $r = 0.93$) than in healthy controls ($R^2 = 0.64$, MAE $= 32.5$ m, $r = 0.80$). The higher POMS $R^2$ is driven by the greater between-subject variation in 6MWD among POMS participants, which gives the model more signal to learn from. The home model gives $R^2 = 0.38$ in POMS and $R^2 = 0.21$ in healthy; MAE is about 10 m larger in POMS (62.6 m) than in healthy (52.1 m). The pattern is consistent with the clinic model: more pathology variability helps the model separate subjects.

\begin{table}[!ht]
\centering
\caption{Per-cohort prediction performance for the two best models. All values are mean [95\% bootstrap CI].}
\label{tab:error}
\small
\begin{tabular}{@{}llcccc@{}}
\toprule
Model & Cohort & $R^2$ & MAE (m) & RMSE (m) & $r$ \\
\midrule
\multirow{2}{*}{\shortstack[l]{Home: Bout+Act-\\Top20+Demo(4)}}
  & POMS ($n=38$)    & 0.38 [0.16, 0.51] & 62.6 [48.3, 78.4] & 77.7 [59.9, 95.9] & 0.63 [0.47, 0.78] \\
  & Healthy ($n=63$) & 0.21 [-0.09, 0.37] & 52.1 [40.5, 66.3] & 72.4 [53.0, 93.3] & 0.47 [0.25, 0.63] \\
\addlinespace
\multirow{2}{*}{\shortstack[l]{Clinic: Gait+CWT+\\WalkSway+Demo}}
  & POMS ($n=38$)    & 0.87 [0.78, 0.92] & 29.0 [22.5, 35.9] & 35.5 [27.9, 43.2] & 0.93 [0.89, 0.96] \\
  & Healthy ($n=63$) & 0.64 [0.42, 0.79] & 32.5 [24.2, 42.5] & 48.8 [31.7, 66.4] & 0.80 [0.70, 0.89] \\
\bottomrule
\end{tabular}
\end{table}

%% ============================================================
\section{Discussion}\label{sec:discuss}
%% ============================================================

We set out to answer whether seven-day home accelerometry can approach the clinic 6MWT for walking-capacity assessment in pediatric-onset MS. The five research questions give a layered answer.

\subsection{What Separates POMS from Healthy Controls (RQ1)}

POMS participants walked with lower cadence, lower intensity, and lower jerk at the clinic, and with shorter, less varied, and less intense bouts at home. The direction is the same in both settings, but the effect sizes are larger at home ($|d|$ up to 1.12 for per-bout jerk) than at the clinic ($|d|$ up to 0.92 for CWT fundamental frequency). This is a useful finding: home monitoring does not just replicate the clinic story, it strengthens it because free-living walking samples more context. The one normalized sway feature that separated groups was ML over ENMO, which was higher in POMS. Many other amplitude-based sway features did not separate the groups because POMS participants both sway less \emph{and} walk less vigorously, so the raw amplitudes shrink together. Normalizing by walking effort exposes the residual instability.

\subsection{What Predicts 6MWD (RQ2)}

Clinic intensity and tempo features were the strongest 6MWD correlates, with $\rho$ up to 0.80. At home, the same features still correlate but at smaller magnitudes ($\rho$ up to 0.37). This gap is expected because the home signal is noisier and captures more non-walking activity. The fact that direction is preserved across settings supports using home features as a proxy for walking capacity, even if single-feature correlations are weaker.

\subsection{What the Clinic and Home Settings Share (RQ3)}

The top finding of RQ3 is that clinic features are most correlated with the 90th-percentile of the home distribution, not with the median or maximum. This makes sense physiologically: a supervised 6MWT asks the subject to walk steadily for six minutes, which samples the upper part of each person's effort distribution. It does not represent the typical home bout, which is shorter and often less vigorous. The correlations were modest ($\rho = 0.28$ to $0.33$) because the two settings still differ in duration, context, and motivation. Even so, the consistent 90th-percentile pattern suggests a practical way to align home and clinic measurements when comparing across cohorts.

\subsection{What Relates to Clinical Scores (RQ4)}

Clinical-score associations separate cleanly between settings. The clinic CWT features correlate with fatigue subscales (MFIS Physical, Cognitive, Psychosocial) in the expected direction: more spectral irregularity, more high-frequency energy, and larger symmetry variability all go with more fatigue. The home features tell a different story that is more about quantity and stability: shorter maximum bout, lower bout-duration variability, slower and more variable stride timing, and fewer activity-to-sedentary transitions all correlate with MFIS and BDI. The two settings therefore measure different aspects of burden. Clinic features capture instability at peak effort; home features capture how often a subject can walk at all and how stable the walking is from bout to bout. Both are clinically useful.

\subsection{How Well Home Predicts 6MWD (RQ5)}

The clinic Ridge model reaches $R^2 = 0.81$ (MAE $= 31.2$ m). The home clinic-free model reaches $R^2 = 0.45$ (MAE $= 56.0$ m). The home model is a partial substitute, not a replacement. Within POMS the clinic model is especially accurate ($R^2 = 0.87$), which is the setting where clinical decisions are most often needed. For adults with MS, prior work reported $r = 0.63$ for free-living 6MWD~\cite{motl2011accelerometry} and $r = 0.71$ for an in-clinic wearable 2MWT estimate~\cite{carpinella2023automatic}. Our POMS home model reaches $r = 0.67$ at a similar performance range without any clinic visit data in the feature pipeline, which is a step forward for remote POMS assessment.

The 25 m MAE gap between clinic and home is consistent across seven different models (Figure~\ref{fig:model_bar}). This tells us the gap is a property of the setting difference (free-living vs.\ supervised), not a property of the model family. Improving home prediction will likely require better data (longer monitoring periods, multi-device sensing, higher-quality bout detection) rather than fancier modeling.

\subsection{Limitations}

First, the cohort is from a single academic center ($n = 101$, POMS $= 38$), which limits generalizability. Second, LOO cross-validation maximizes training data but may give slightly optimistic estimates compared to a held-out test set. Third, the study is cross-sectional; longitudinal follow-up is needed to show whether the home-predicted 6MWD tracks meaningful changes over time. Fourth, home monitoring depends on wear compliance and on walking-bout detection accuracy, both of which add noise. Fifth, the home model MAE of 56.0 m is wider than the minimal clinically important difference for the 6MWT in adult MS (25 to 30 m), so home predictions should complement, not replace, clinic testing.

\subsection{Future Directions}

Future work should extend this pipeline to longitudinal POMS monitoring to test whether home features can detect meaningful decline before a clinic visit would. Larger and more diverse POMS cohorts across multiple sites would improve generalizability. Combining accelerometry with other passive sensors (smartphone touches, sleep, heart rate) may narrow the clinic-to-home gap. Finally, the 90th-percentile finding of RQ3 suggests that home-derived peak effort features are the right anchor for aligning free-living data with supervised clinic tests, which could be formalized into a calibration procedure.

%% ============================================================
\section{Conclusion}\label{sec:conclusion}
%% ============================================================

We used hip-worn accelerometry in the clinic and at home to study gait in 101 participants (38 with pediatric-onset MS and 63 healthy controls). POMS participants walked with lower cadence, lower intensity, and higher relative lateral sway than controls; at home they showed shorter and less varied walking bouts. Clinic gait intensity features tracked 6MWD most strongly, while home features carried the same signal at smaller magnitudes. Clinic and home measurements were most correlated when the home value was taken at the 90th percentile of a subject's walking bouts, meaning the supervised 6MWT samples near-peak free-living effort. CWT spectral features at the clinic and stride-timing variability at home correlated with fatigue scores. Ridge regression predicted 6MWD with $R^2 = 0.81$ (MAE $= 31.2$ m) in the clinic and $R^2 = 0.45$ (MAE $= 56.0$ m) in a fully clinic-free home pipeline. The home model is a partial substitute for the clinic 6MWT and enables remote monitoring between clinic visits, while the clinic model remains the best option when precise 6MWD estimation is needed.

%% ============================================================
\bibliographystyle{unsrt}
\bibliography{references}

%% ============================================================
\clearpage
\appendix
\section*{Supplementary Figures}\label{sec:supp}
%% ============================================================

\begin{figure}[H]
    \centering
    \includegraphics[width=0.95\linewidth,height=0.32\textheight,keepaspectratio]{figures/fig_supp_wear_median_bout.png}
    \caption{Supplementary Figure S1. (a) Device wear time over the seven-day home monitoring period, shown as violin plots with jittered subject points for POMS and healthy controls. Median wear time is 59 hours for POMS and 53 hours for healthy controls, with no significant group difference (Mann-Whitney $U$, $p = 0.19$). (b) Per-subject median bout duration by cohort (one data point per subject). Medians are 14 s for both groups, with POMS showing a wider lower tail. (c) Per-subject scatter of number of walking bouts against median bout duration. Point size is scaled by 6MWD. Subjects with more bouts and longer median bout durations tend to have longer 6MWD. Axes use log scales to accommodate the wide range of bout counts and durations.}
    \label{fig:supp_wear_median_bout}
\end{figure}

\vspace{-0.4em}

\begin{figure}[H]
    \centering
    \includegraphics[width=0.95\linewidth,height=0.28\textheight,keepaspectratio]{figures/bout_duration_boxplots_per_subject.png}
    \caption{Supplementary Figure S2. Per-subject bout-duration distribution. Each box shows one subject, sorted by median bout duration. POMS subjects are in the top row (red) and healthy subjects are in the bottom row (blue). The visualization shows the within-subject spread of bout durations and how the spread varies between groups.}
    \label{fig:supp_boxplots}
\end{figure}

\vspace{-0.4em}

\begin{figure}[H]
    \centering
    \includegraphics[width=0.80\linewidth,height=0.28\textheight,keepaspectratio]{figures/fig_supp_bout_threshold_sensitivity.png}
    \caption{Supplementary Figure S3. Sensitivity of home 6MWD prediction MAE to the minimum walking-bout duration threshold. At each threshold (10, 30, 60, 120, 240 s) the cached per-bout features are filtered, re-aggregated, combined with the 29 whole-recording activity features, and fed to the headline Bout+Act-Top20+Demo pipeline (Ridge $\alpha=20$, leave-one-out cross-validation). Error bars are 95\% percentile bootstrap confidence intervals (B = 2000). The 10\,s threshold yields the best MAE of 56.0\,m; stricter thresholds reduce the number of usable bouts (mean bouts per subject drop from 1{,}454 at 10\,s to 11 at 240\,s) and at 240\,s leave 9 of 101 subjects with zero bouts, degrading MAE to 65.5\,m.}
    \label{fig:supp_bout_threshold}
\end{figure}

\end{document}
